% GENERATED FILE. Edit canonical lessons, then run npm run generate:books.
% canonical-source-hash: fnv1a64:f79e13e165a6881e
% canonical-lessons: TE-C494-digulu, TE-C494-cintincu, TE-C494-kalavarapadu, TE-C494-kangarupadu, TE-C494-atruta

\chapter{Worry and Unease}
\label{ch:te-worry-and-unease}
\hlchaptermodality{494}

\begin{chapteropening}
\textbf{By the end of this chapter:} \emph{I can say that I feel low, worried, upset or flustered.}

Complete the last lesson of chapter 494: I can say that I feel low, worried, upset or flustered.
\end{chapteropening}

\section[\texorpdfstring{\te{దిగులు} (digulu)}{దిగులు (digulu)}]{\te{దిగులు} (digulu) — gloom; low spirits; worry}

\label{lesson:TE-C494-digulu}



\begin{quote}
Before the new one: say the Telugu for a watchman, then the Telugu for a milkman.
\end{quote}

\subsection*{You'll want to know: \te{దిగులు}}
\textbf{\te{దిగులు}} — \emph{digulu} — \textquotedblleft{}gloom; low spirits; worry\textquotedblright{}.

\begin{quote}
\te{అమ్మ} \te{దూరంగా} \te{ఉంటే} \te{నాకు} \te{దిగులుగా} \te{ఉంటుంది}.
\end{quote}

\begin{quote}
\emph{amma dūraṅgā uṇṭē nāku digulugā uṇṭundi.}
\end{quote}

\begin{quote}
When mum is far away, I feel low.
\end{quote}

Said with the dative: \te{నాకు} \te{దిగులుగా} \te{ఉంది}, I feel low.

\begin{cousinweb}
\textbf{\te{దిగులు}} (\emph{digulu}) is native Dravidian: DEDR 3202 lists Telugu \emph{digulu}, \textquotedblleft{}fear, alarm\textquotedblright{}, beside Tamil \emph{tikil}, \textquotedblleft{}fright, terror\textquotedblright{}. In everyday Telugu the fear has softened into a heavy, anxious sadness.
\end{cousinweb}

\begin{grammarlens}[title={What you've built}]
One more word for describing what you lived and felt.
\end{grammarlens}

\subsection*{Your turn}
\begin{itemize}
\raggedright
  \item \emph{Say it:} \emph{digulu}
  \item \emph{Say it:} \emph{digulu}, once more
  \item \emph{Say it:} \emph{amma dūraṅgā uṇṭē nāku digulugā uṇṭundi.} — When mum is far away, I feel low.
  \item \emph{Recall:} say the Telugu for a watchman, then the Telugu for a milkman, then say \emph{digulu} again
\end{itemize}

\begin{tcolorbox}[breakable,colback=teal!4,colframe=teal!35!black,title={Before you move on}]
What does \te{దిగులు} mean? (\textquotedblleft{}Gloom; low spirits; worry\textquotedblright{}.) Say it once more.
\end{tcolorbox}
\section[\texorpdfstring{\te{చింతించు} (cintiñcu)}{చింతించు (cintiñcu)}]{\te{చింతించు} (cintiñcu) — to worry; to regret}

\label{lesson:TE-C494-cintincu}



\begin{quote}
Before the new one: say the Telugu for a milkman, then the Telugu for gloom.
\end{quote}

\subsection*{You'll want to know: \te{చింతించు}}
\textbf{\te{చింతించు}} — \emph{cintiñcu} — \textquotedblleft{}to worry; to regret\textquotedblright{}.

\begin{quote}
\te{నాన్న} \te{డబ్బు} \te{గురించి} \te{చాలా} \te{చింతించాడు}.
\end{quote}

\begin{quote}
\emph{nānna ḍabbu guriñci cālā cintiñcāḍu.}
\end{quote}

\begin{quote}
Dad worried a lot about money.
\end{quote}

Past (I …): \te{చింతించాను}, \emph{cintiñcānu}.

What you worry about takes \te{గురించి}. The noun is \te{చింత}, worry or care.

\begin{cousinweb}
\textbf{\te{చింతించు}} (\emph{cintiñcu}) is borrowed from Sanskrit \emph{cintayati}, \textquotedblleft{}thinks; cares for\textquotedblright{} (CDIAL 4815), and takes the Telugu ending -\te{ఇంచు}, which turns a Sanskrit stem into a verb. Thinking too much about a thing becomes worrying about it.
\end{cousinweb}

\begin{grammarlens}[title={What you've built}]
One more word for describing what you lived and felt.
\end{grammarlens}

\subsection*{Your turn}
\begin{itemize}
\raggedright
  \item \emph{Say it:} \emph{cintiñcu}
  \item \emph{Say it:} \emph{cintiñcu}, once more
  \item \emph{Say it:} \emph{nānna ḍabbu guriñci cālā cintiñcāḍu.} — Dad worried a lot about money.
  \item \emph{Recall:} say the Telugu for a milkman, then the Telugu for gloom, then say \emph{cintiñcu} again
\end{itemize}

\begin{tcolorbox}[breakable,colback=teal!4,colframe=teal!35!black,title={Before you move on}]
What does \te{చింతించు} mean? (\textquotedblleft{}To worry; to regret\textquotedblright{}.) Say it once more.
\end{tcolorbox}
\section[\texorpdfstring{\te{కలవరపడు} (kalavarapaḍu)}{కలవరపడు (kalavarapaḍu)}]{\te{కలవరపడు} (kalavarapaḍu) — to be agitated; to be upset}

\label{lesson:TE-C494-kalavarapadu}



\begin{quote}
Before the new one: say the Telugu for gloom, then the Telugu for to worry.
\end{quote}

\subsection*{You'll want to know: \te{కలవరపడు}}
\textbf{\te{కలవరపడు}} — \emph{kalavarapaḍu} — \textquotedblleft{}to be agitated; to be upset\textquotedblright{}.

\begin{quote}
\te{ఆ} \te{వార్త} \te{విని} \te{అమ్మ} \te{కలవరపడింది}.
\end{quote}

\begin{quote}
\emph{ā vārta vini amma kalavarapaḍindi.}
\end{quote}

\begin{quote}
Hearing that news, mum was upset.
\end{quote}

Past (I …): \te{కలవరపడ్డాను}, \emph{kalavarapaḍḍānu}.

The noun is \te{కలవరం}, confusion of mind. A noun of feeling + -\te{పడు} means to feel it.

\begin{cousinweb}
\textbf{\te{కలవరపడు}} (\emph{kalavarapaḍu}) is native Dravidian: DEDR 1306 lists Telugu \emph{kalavaramu}, \textquotedblleft{}confusion, anxiety\textquotedblright{}, and \emph{kalavara-paḍu}, \textquotedblleft{}to be confused\textquotedblright{}, beside Tamil \emph{kalavaram}, \textquotedblleft{}perturbation\textquotedblright{}. The second part is the native verb \emph{paḍu}, \textquotedblleft{}to feel, to suffer\textquotedblright{} (DEDR 3853).
\end{cousinweb}

\begin{grammarlens}[title={What you've built}]
One more word for describing what you lived and felt.
\end{grammarlens}

\subsection*{Your turn}
\begin{itemize}
\raggedright
  \item \emph{Say it:} \emph{kalavarapaḍu}
  \item \emph{Say it:} \emph{kalavarapaḍu}, once more
  \item \emph{Say it:} \emph{ā vārta vini amma kalavarapaḍindi.} — Hearing that news, mum was upset.
  \item \emph{Recall:} say the Telugu for gloom, then the Telugu for to worry, then say \emph{kalavarapaḍu} again
\end{itemize}

\begin{tcolorbox}[breakable,colback=teal!4,colframe=teal!35!black,title={Before you move on}]
What does \te{కలవరపడు} mean? (\textquotedblleft{}To be agitated; to be upset\textquotedblright{}.) Say it once more.
\end{tcolorbox}
\section[\texorpdfstring{\te{కంగారుపడు} (kaṅgārupaḍu)}{కంగారుపడు (kaṅgārupaḍu)}]{\te{కంగారుపడు} (kaṅgārupaḍu) — to panic; to get flustered}

\label{lesson:TE-C494-kangarupadu}



\begin{quote}
Before the new one: say the Telugu for to worry, then the Telugu for to be agitated.
\end{quote}

\subsection*{You'll want to know: \te{కంగారుపడు}}
\textbf{\te{కంగారుపడు}} — \emph{kaṅgārupaḍu} — \textquotedblleft{}to panic; to get flustered\textquotedblright{}.

\begin{quote}
\te{రైలు} \te{రాకపోతే} \te{అతను} \te{కంగారుపడతాడు}.
\end{quote}

\begin{quote}
\emph{railu rākapōtē atanu kaṅgārupaḍatāḍu.}
\end{quote}

\begin{quote}
If the train does not come, he gets flustered.
\end{quote}

Past (I …): \te{కంగారుపడ్డాను}, \emph{kaṅgārupaḍḍānu}.

The noun \te{కంగారు} is fluster or panic: \te{కంగారుగా}, in a flurry.

\begin{cousinweb}
\textbf{\te{కంగారుపడు}} (\emph{kaṅgārupaḍu}) is built from \emph{kaṅgāru}, \textquotedblleft{}fluster, panic\textquotedblright{}, and the native verb \emph{paḍu}, \textquotedblleft{}to feel, to suffer\textquotedblright{} (DEDR 3853). The origin of \emph{kaṅgāru} itself is not settled.
\end{cousinweb}

\begin{grammarlens}[title={What you've built}]
One more word for describing what you lived and felt.
\end{grammarlens}

\subsection*{Your turn}
\begin{itemize}
\raggedright
  \item \emph{Say it:} \emph{kaṅgārupaḍu}
  \item \emph{Say it:} \emph{kaṅgārupaḍu}, once more
  \item \emph{Say it:} \emph{railu rākapōtē atanu kaṅgārupaḍatāḍu.} — If the train does not come, he gets flustered.
  \item \emph{Recall:} say the Telugu for to worry, then the Telugu for to be agitated, then say \emph{kaṅgārupaḍu} again
\end{itemize}

\begin{tcolorbox}[breakable,colback=teal!4,colframe=teal!35!black,title={Before you move on}]
What does \te{కంగారుపడు} mean? (\textquotedblleft{}To panic; to get flustered\textquotedblright{}.) Say it once more.
\end{tcolorbox}
\section[\texorpdfstring{\te{ఆత్రుత} (ātruta)}{ఆత్రుత (ātruta)}]{\te{ఆత్రుత} (ātruta) — eagerness; impatience}

\label{lesson:TE-C494-atruta}



\begin{quote}
Before the new one: say the Telugu for to be agitated, then the Telugu for to panic.
\end{quote}

\subsection*{You'll want to know: \te{ఆత్రుత}}
\textbf{\te{ఆత్రుత}} — \emph{ātruta} — \textquotedblleft{}eagerness; impatience\textquotedblright{}.

\begin{quote}
\te{ఫలితాల} \te{కోసం} \te{విద్యార్థులు} \te{ఆత్రుతగా} \te{ఎదురుచూశారు}.
\end{quote}

\begin{quote}
\emph{phalitāla kōsaṁ vidyārthulu ātrutagā edurucūśāru.}
\end{quote}

\begin{quote}
The students waited eagerly for the results.
\end{quote}

Mostly as an adverb: \te{ఆత్రుతగా}, eagerly, impatiently.

\begin{cousinweb}
\textbf{\te{ఆత్రుత}} (\emph{ātruta}) is usually traced to Sanskrit \emph{ātura}, \textquotedblleft{}suffering, troubled in body or mind\textquotedblright{} (Monier-Williams), with an abstract ending. The restless feeling became eager impatience.
\end{cousinweb}

\begin{grammarlens}[title={What you've built}]
One more word for describing what you lived and felt.
\end{grammarlens}

\subsection*{Your turn}
\begin{itemize}
\raggedright
  \item \emph{Say it:} \emph{ātruta}
  \item \emph{Say it:} \emph{ātruta}, once more
  \item \emph{Say it:} \emph{phalitāla kōsaṁ vidyārthulu ātrutagā edurucūśāru.} — The students waited eagerly for the results.
  \item \emph{Recall:} say the Telugu for to be agitated, then the Telugu for to panic, then say \emph{ātruta} again
\end{itemize}

\begin{tcolorbox}[breakable,colback=teal!4,colframe=teal!35!black,title={Before you move on}]
What does \te{ఆత్రుత} mean? (\textquotedblleft{}Eagerness; impatience\textquotedblright{}.) Say it once more.
\end{tcolorbox}
